\chapter{Fonctions polynomiales et rationnelles}
\label{workbook:chapter:05}
\section{Architecture des fonctions polynomiales et rationnelles}
\label{workbook:section:05.01}
\begin{exerciseblock}{Architecture des fonctions polynomiales et rationnelles}
\item Pour $P(x)=-2x^4+3x^2-1$, donner degré, coefficient dominant, parité et comportement aux deux infinis.
\item Pour $R(x)=\dfrac{x^2-4}{x^2-x-6}$, déterminer domaine, factorisation et éventuelles simplifications.
\item Comparer les informations fournies par numérateur et dénominateur d'une fonction rationnelle.
\end{exerciseblock}
\section{Prévoir un graphe à partir de l'expression}
\label{workbook:section:05.02}
\begin{exerciseblock}{Prévoir un graphe à partir de l'expression}
\item Sans tracer précisément, esquisser les grandes caractéristiques de $f(x)=-(x+2)^2(x-1)^3$.
\item Prévoir les asymptotes et les intersections principales de $g(x)=\dfrac{2x+1}{x-3}$.
\item Construire un polynôme de degré $5$ ayant une racine double en $-1$, une racine triple en $2$ et coefficient dominant positif.
\end{exerciseblock}
\section{Comportement local et comportement global}
\label{workbook:section:05.03}
\begin{exerciseblock}{Comportement local et comportement global}
\item Comparer le comportement près de $x=1$ de $(x-1)$, $(x-1)^2$ et $(x-1)^3$.
\item Étudier les limites qualitatives de $\dfrac{1}{x-2}$ lorsque $x\to2^-$, $x\to2^+$ et $|x|\to\infty$.
\item Pour $P(x)=x^5-100x^2+7$, expliquer pourquoi le terme dominant gouverne le comportement global mais pas nécessairement le graphe près de l'origine.
\end{exerciseblock}
\section{Du facteur au graphe}
\label{workbook:section:05.04}
\begin{exerciseblock}{Du facteur au graphe}
\item Construire le tableau de signes de $P(x)=(x+3)(x-1)^2(x-4)$.
\item Effectuer la division de $x^2+1$ par $x-1$ et en déduire l'asymptote oblique de $\dfrac{x^2+1}{x-1}$.
\item Faire une étude qualitative complète de $R(x)=\dfrac{x+1}{x^2-4}$ : domaine, zéros, signes et asymptotes.
\end{exerciseblock}
\section*{Synthèse du chapitre}
\begin{synthesisbox}
\item Étudier et esquisser $f(x)=\dfrac{x^2-1}{x^2-4}$ en indiquant toutes les caractéristiques utiles.
\item Construire un polynôme minimal satisfaisant : racines $-2$ et $3$, la première double, la seconde simple, et $P(0)=-12$.
\item Donner deux fonctions rationnelles ayant la même asymptote horizontale mais des domaines et des zéros différents.
\end{synthesisbox}
